\documentclass[UTF8,a4paper,11pt]{ctexart}
\usepackage[a4paper,margin=2.3cm]{geometry}
\usepackage{amsmath,amssymb,mathtools}
\usepackage{enumitem}
\usepackage[hidelinks]{hyperref}
\usepackage{fancyhdr}
\setlist[itemize]{leftmargin=2.2em,itemsep=0.25em,topsep=0.35em}
\setlength{\parindent}{2em}\setlength{\parskip}{0.35em}
\pagestyle{fancy}\fancyhf{}\fancyhead[L]{离散数学与结构（I）}\fancyhead[R]{Lec 19}\fancyfoot[C]{\thepage}\setlength{\headheight}{14pt}
\title{无损压缩、Huffman 与侧信息}\author{zhangjunshuo}\date{2026 年 11 月}
\begin{document}
\maketitle
\begin{quote}
本讲根据 2026 年 11 月 19 日的三页手写稿和两段课堂转写整理。上一份笔记是 Lec18（11 月 17 日），本次按顺序记为 \textbf{Lec19}。手写稿是公式与结构的主要依据，转写用于补课堂解释。

本讲主线：无损变长码 $\to$ Kraft/Huffman $\to$ 几乎无损定长码（典型集）$\to$ 侧信息与 Slepian--Wolf。

每个分节开头都先交代本节用到的记号；所有符号在首次使用处给出定义。
\end{quote}
\setcounter{tocdepth}{1}
\tableofcontents\clearpage

\section{无损压缩与变长码}


这一讲进入\textbf{无损压缩}：把信息源 $X\sim P$ 编码成二进制串，要求解码后完全恢复原文，并让平均码长尽量短。核心结论是：\textbf{熵是平均码长的下界，Huffman 码在前缀码中最优，且最优平均码长不超过 $H+1$。}

\textbf{本节用到的记号：} $\Omega$ 为信源字母表；$P$ 为分布；$f:\Omega\to\{0,1\}^*$ 为编码；$g:\{0,1\}^*\to\Omega$ 为解码；$l(x)=|f(x)|$ 为码长；$L=E[l(X)]$ 为平均码长；$H(X)$ 为熵。

\subsection{问题设定}


无损压缩由一对函数组成：

\[
f:\Omega\to\{0,1\}^*,\qquad g:\{0,1\}^*\to\Omega,
\]

满足

\[
g(f(x))=x\qquad\forall x\in\Omega.
\]

也就是说，$f$ 必须是单射；否则两个不同信源符号压成同一个串，就无法解码。

\textbf{优化目标.} 最小化平均码长

\[
L(P,f)=E_{X\sim P}\bigl[|f(X)|\bigr]
=\sum_{x\in\Omega}P(x)|f(x)|.
\]

\textbf{直观.} 出现概率大的符号用短码，出现概率小的符号用长码。 Morse 码就是这样：E 用一点，不常用的字母用较长码。

\subsection{变长码与即时解码}


只要求“最终能解码”还不够：实际传输是流式的，收到一个码字后，最好立刻知道它对应哪个符号，而不必等后续比特。

\textbf{前缀码（prefix-free code）.} 没有任何码字是另一个码字的前缀。例如

\[
\{0,10,11\}
\]

是前缀码；而 $\{0,01\}$ 不是，因为 \texttt{0} 是 \texttt{01} 的前缀。

前缀码可以画成一棵二叉树的叶子：每个叶子对应一个码字，从根到叶子的路径是码字。前缀条件保证任何一个码字都不会停在另一个码字的路径中间。

\textbf{注意.} Morse 码本身不是前缀码；它靠字母之间的间隔、单词之间的更长间隔来区分。本课程主要讨论前缀码。

\subsection{平均码长的下界：熵}


对任意前缀码，平均码长满足

\[
\boxed{
E[|f(X)|]\ge H(X)
}.
\]

\textbf{证明（用 Kraft + Jensen）.} 前缀码的码长 $l(x)$ 满足 Kraft 不等式

\[
\sum_{x}2^{-l(x)}\le1.
\]

令

\[
Q(x)=\frac{2^{-l(x)}}{Z},\qquad
Z=\sum_x2^{-l(x)}\le1.
\]

则 $Q$ 是一个分布。计算 $P$ 对 $Q$ 的 KL 散度：

\[
\begin{aligned}
D(P\Vert Q)
&=\sum_xP(x)\log\frac{P(x)}{Q(x)}\\
&=\sum_xP(x)\log\frac{P(x)}{2^{-l(x)}/Z}\\
&=\sum_xP(x)l(x)-H(P)+\log Z\\
&=E[L]-H(P)+\log Z.
\end{aligned}
\]

由 KL 非负性，$D(P\Vert Q)\ge0$。又 $Z\le1$，所以 $\log Z\le0$，于是

\[
E[L]-H(P)
=D(P\Vert Q)-\log Z
\ge -\log Z\ge0,
\]

故

\[
E[L]\ge H(P).
\]

等号需要 $P=Q$ 且 $Z=1$，即 $P(x)=2^{-l(x)}$，码长正好是 $-\log P(x)$（一般为整数时才可能完全取等）。

\subsection{上界：Shannon 码}


取码长

\[
l(x)=\left\lceil\log_2\frac1{P(x)}\right\rceil.
\]

则

\[
2^{-l(x)}\le P(x),
\]

所以

\[
\sum_x2^{-l(x)}\le\sum_xP(x)=1,
\]

满足 Kraft 不等式，因此存在前缀码具有这些码长。其平均长度满足

\[
l(x)<\log_2\frac1{P(x)}+1,
\]

所以

\[
E[L]<H(X)+1.
\]

结合下界：

\[
\boxed{
H(X)\le E[L]<H(X)+1
}.
\]

也就是说，用前缀码最多比熵多花不到 $1$ bit。

\subsection{分块编码}


把 $n$ 个独立同分布符号一起编码，码长记为 $L_n$，则

\[
nH(X)\le E[L_n]<nH(X)+1.
\]

平均到每个符号：

\[
H(X)\le\frac{E[L_n]}{n}<H(X)+\frac1n.
\]

当 $n\to\infty$ 时，每符号平均码长趋于熵。这是“熵是压缩极限”的最基本表述。

\subsection{常见错误}


\begin{itemize}
\item 把“能解码”与“前缀码”混淆：前缀码是更强的即时解码条件。
\item 认为熵是“某个具体编码的长度”；熵是平均码长的下界，不是每个码字的长度。
\item 把 Shannon 码的 $+1$ 漏掉；上界是 $H+1$，不是 $H$。
\item 对非前缀码直接套 Kraft 不等式；Kraft 是前缀码/前缀码树的条件（唯一可译码也有类似但更复杂的界）。
\end{itemize}

\section{Kraft 不等式与 Huffman 编码}


上一节给出了前缀码平均码长的上下界；这一节说明“哪些码长组合可以实现”，并给出达到最优的构造算法：\textbf{Huffman 编码}。

\textbf{本节用到的记号：} 码长 $l_1,\dots,l_m$；Kraft 和 $K=\sum_i2^{-l_i}$；$H(P)$ 为熵；$L^*$ 为最优前缀码的平均码长。

\subsection{Kraft 不等式}


\textbf{定理.} 存在以 $l_1,\dots,l_m$ 为码长的前缀码，当且仅当

\[
\boxed{
\sum_{i=1}^m2^{-l_i}\le1
}.
\]

\textbf{必要性证明（随机无限串）.} 独立均匀地生成一个无限比特串 $S\in\{0,1\}^{\mathbb N}$。对每个码字 $c_i$，事件“$c_i$ 是 $S$ 的前缀”的概率是

\[
2^{-l_i}.
\]

前缀码保证这些事件两两不交（一个无限串不可能同时以两个不同码字为前缀）。因此

\[
\sum_{i=1}^m2^{-l_i}
=\sum_i\Pr[c_i\text{ 是 }S\text{ 的前缀}]
\le1.
\]

\textbf{充分性（构造树）.} 把长度 $l_i$ 的码字看成深度为 $l_i$ 的叶子。按长度从小到大放置：每个长度为 $l_i$ 的叶子占据一个“容量为 $2^{-l_i}$”的区间；Kraft 和 $\le1$ 保证这些区间不相交，于是可以构造出一棵前缀码树。

\textbf{直观.} Kraft 和就是所有码字叶子在无限比特串空间中所占的“总测度”；前缀码的叶子不能重叠，所以总测度不超过 $1$。

\subsection{Huffman 编码}


\textbf{算法.}

1. 把每个符号 $x_i$ 看成一个权为 $P(x_i)$ 的叶子；
2. 每次选出权最小的两个结点，合并成一个新结点，权为两者之和；
3. 重复直到只剩一个根；
4. 从根到叶子的路径上，左分支记 \texttt{0}、右分支记 \texttt{1}，得到每个符号的码字。

\textbf{性质.}

\begin{itemize}
\item Huffman 码是\textbf{前缀码}；
\item 在所有前缀码中，Huffman 码的\textbf{平均码长最小}；
\item 最优平均码长满足
\end{itemize}

\[
  H(P)\le L^*<H(P)+1.
\]

\textbf{最优性证明思路（交换论证）.} 在最优前缀码树中，概率最小的两个符号一定可以放在最深层的兄弟位置：如果它们不在最深处，把它们与最深处的叶子交换，平均码长不会增加。于是可以把这两个最小概率合并成一个“超符号”，递归构造；这正是 Huffman 算法。由此归纳可得最优性。

\textbf{例子（思路）.} 若概率为

\[
0.4,\ 0.2,\ 0.2,\ 0.1,\ 0.1,
\]

则先合并两个 $0.1$ 得到 $0.2$，再合并两个 $0.2$ 得到 $0.4$，最后合并 $0.2$ 与 $0.4$；一种最优码长是 $2,2,2,3,3$（不同合并顺序可能得到不同但等价的树）。概率大的符号码短，概率小的码长。

\subsection{变长码的极限}


对 $n$ 个独立同分布符号一起做前缀编码，最优平均码长 $L_n^*$ 满足

\[
nH(P)\le L_n^*<nH(P)+1.
\]

因此每符号平均码长

\[
\frac{L_n^*}{n}\to H(P).
\]

这就是“熵是可逆压缩的极限速率”。

\subsection{常见错误}


\begin{itemize}
\item 把 Kraft 不等式写成“$\sum2^{-l_i}=1$”；它是不等式，$\le1$ 即可。
\item 认为 Huffman 对\textbf{任意}唯一可译码最优；它是在\textbf{前缀码}类中最优。前缀码类已经能达到唯一可译码的最优渐近速率。
\item 把“最短码”给最不常出现的符号；恰好相反，最短码给概率最大的符号。
\item 忘记 Huffman 需要已知分布；分布未知时要用通用压缩/随机压缩（下一节）。
\end{itemize}

\section{定长压缩与典型集}


变长码能做到 $H(X)\le E[L]<H(X)+1$，但每条消息长度不同。很多实际场景要求\textbf{定长编码}：把 $n$ 个符号压成固定 $L$ 个比特。代价是允许一个很小的错误概率，这叫\textbf{几乎无损压缩}（almost lossless compression）。

\textbf{本节用到的记号：} $X^n=(X_1,\dots,X_n)$ iid $\sim P$；$P^n(x)=\prod_iP(x_i)$；$f:\Omega^n\to\{0,1\}^L$ 为定长编码；$g$ 为解码；$T_\varepsilon^{(n)}$ 为典型集；$H=H(P)$。

\subsection{精确无损定长的限制}


如果要求对\textbf{所有} $x\in\Omega^n$ 都精确解码，那么需要

\[
2^L\ge|\Omega|^n
\quad\Longrightarrow\quad
L\ge n\log|\Omega|.
\]

这通常远大于 $nH(P)$；所以精确无损定长压缩无法达到熵极限。要接近熵，必须允许少量错误。

\subsection{几乎无损定长压缩}


允许错误概率 $\varepsilon$：要求

\[
\Pr[g(f(X^n))\ne X^n]\le\varepsilon.
\]

可以构造 $L\approx nH(P)$ 的编码，使错误概率随 $n$ 指数小。

\textbf{核心工具：典型集.} 定义

\[
T_\varepsilon^{(n)}
=\left\{x\in\Omega^n:
\left|-\frac1n\log P^n(x)-H\right|\le\varepsilon
\right\}.
\]

典型集有两个关键性质：

1. \textbf{概率大}：由大数定律，
\[
   \Pr[X^n\in T_\varepsilon^{(n)}]\to1
   \quad(n\to\infty).
\]
2. \textbf{数量少}：每个典型序列的概率至少约 $2^{-n(H+\varepsilon)}$，而总概率不超过 $1$，所以
\[
   |T_\varepsilon^{(n)}|\le2^{n(H+\varepsilon)}.
\]

\textbf{编码构造.} 取

\[
L=\left\lceil n(H+\varepsilon)\right\rceil.
\]

把典型集中的序列任意编号为 $1,\dots,2^L$；非典型序列随便编码（解码时会错，但概率小）。解码时按编号恢复。

\textbf{错误概率.} 只有非典型序列会出错，因此

\[
\Pr[g(f(X^n))\ne X^n]
\le\Pr[X^n\notin T_\varepsilon^{(n)}]
\le\varepsilon
\]

（对足够大的 $n$）。更精细的集中不等式可给出 $e^{-\Theta(n\varepsilon^2)}$ 级的错误概率。

\subsection{随机压缩：分布未知时怎么办}


上面的构造需要知道 $P$ 才能定义典型集。若分布未知，可以用\textbf{随机压缩}：

\begin{itemize}
\item 随机选一个函数 $f:\Omega^n\to\{0,1\}^L$；
\item 解码时，给定消息 $s$，在 $f^{-1}(s)$ 中选概率最大的 $x$：
\end{itemize}

\[
  g(s)=\arg\max_{x\in f^{-1}(s)}P^n(x).
\]
\begin{itemize}
\item 若 $L>nH(P)+\varepsilon$，则对每个高概率的 $x$，它与其它高概率点碰撞的概率很小，因此可以高概率正确解码。
\end{itemize}


\textbf{随机线性压缩.} 也可以取 $f$ 为随机矩阵（在 $\mathbb F_2$ 上）作用在 $x$ 上；两个不同向量碰撞的概率仍约为输出空间大小的倒数，因此同样的结论成立。这样连分布的具体形式都不需要知道，只需要知道熵的上界。

\subsection{固定长度 vs 变长码}


\begin{itemize}
\item \textbf{变长码}：每条消息长度可变；精确无损可做到 $E[L]\approx H$；也可分块做几乎无损；工具是 Kraft/Huffman。
\item \textbf{定长码}：每条消息固定 $L$ bit；精确无损需要 $L\ge n\log|\Omega|$；几乎无损可用 $L\approx nH$，错误概率指数小；工具是典型集/大数定律。
\end{itemize}


\subsection{常见错误}


\begin{itemize}
\item 把“几乎无损”当成“精确无损”；定长编码在 $L<n\log|\Omega|$ 时必然有错。
\item 把典型集大小写成 $2^{nH}$ 而漏掉 $\varepsilon$；应是 $2^{n(H\pm\varepsilon)}$。
\item 认为随机压缩需要知道分布；随机函数/随机矩阵的构造不需要完整分布，只需熵的界。
\item 把错误概率的指数写成 $e^{-n\varepsilon}$；典型集的偏差通常是 $e^{-n\varepsilon^2}$ 量级。
\end{itemize}

\section{侧信息与 Slepian–Wolf 压缩}


这一节讨论\textbf{分布式无损压缩}：编码器只能看到 $X$，但解码器还能看到与 $X$ 相关的旁信息 $Y$。问题是：既然解码器已经有 $Y$，$X$ 还需要多少比特？

\textbf{本节用到的记号：} $(X^n,Y^n)$ iid $\sim P_{X,Y}$；$H(X\mid Y)$ 为条件熵；$L_1,L_2$ 为两个编码器的码率（每符号比特数）；$\varepsilon$ 为允许的错误概率。

\subsection{只有解码器有侧信息}


\textbf{设定.}

\begin{itemize}
\item 编码器看到 $X^n$，发送 $L_1$ bit/符号的消息；
\item 解码器看到消息和 $Y^n$；
\item 要求以高概率恢复 $X^n$。
\end{itemize}


\textbf{结论.} 只要

\[
L_1>H(X\mid Y),
\]

就存在这样的压缩方案（$n$ 充分大、错误概率任意小）。

\textbf{直观.} 解码器已经知道 $Y$，所以 $X$ 剩下的不确定性只有 $H(X\mid Y)$；编码器只需提供这些“剩余信息”。

\textbf{构造（随机 binning）.}

1. 把 $X^n$ 的所有可能序列随机分到 $2^{nL_1}$ 个桶（bin）中；
2. 编码器只发送桶号；
3. 解码器收到桶号后，在桶中寻找与 $Y^n$ 联合典型的 $X^n$，输出它。

若 $L_1>H(X\mid Y)$，则每个桶中与给定 $Y^n$ 联合典型的 $X^n$ 大约只有

\[
2^{nH(X\mid Y)}/2^{nL_1}<1
\]

个，因此解码唯一、正确概率高。

\textbf{错误概率.} 用联合典型性/集中不等式可证错误概率为 $e^{-\Theta(n)}$ 级（$n\to\infty$）。

\subsection{两个编码器分别压缩（Slepian–Wolf）}


\textbf{设定.}

\begin{itemize}
\item 编码器 1 看 $X^n$，发送速率 $L_1$；
\item 编码器 2 看 $Y^n$，发送速率 $L_2$；
\item 解码器同时收到两个消息，要恢复 $(X^n,Y^n)$。
\end{itemize}


\textbf{Slepian–Wolf 区域.} 可达速率对 $(L_1,L_2)$ 满足

\[
\boxed{
L_1\ge H(X\mid Y),\qquad
L_2\ge H(Y\mid X),\qquad
L_1+L_2\ge H(X,Y)
}
\]

（严格地说要加任意小 $\varepsilon$）。

\textbf{三个下界的直观.}

\begin{itemize}
\item $L_1\ge H(X\mid Y)$：即使解码器有完整的 $Y$，$X$ 仍有 $H(X\mid Y)$ 的不确定性；
\item $L_2\ge H(Y\mid X)$：对称；
\item $L_1+L_2\ge H(X,Y)$：两个消息合起来的总信息量不能少于联合熵。
\end{itemize}


\textbf{可达性.} 用随机 binning：把 $X^n$ 分到 $2^{nL_1}$ 个桶、$Y^n$ 分到 $2^{nL_2}$ 个桶；解码器在“桶对”中找联合典型的 $(X^n,Y^n)$。只要速率在区域内部，联合典型对唯一，错误概率指数小。

\subsection{与单源压缩的关系}


\begin{itemize}
\item 若 $Y$ 是常数（没有侧信息），则 $H(X\mid Y)=H(X)$，回到普通无损压缩 $L\ge H(X)$；
\item 若 $Y=X$，则 $H(X\mid Y)=0$，不需要传 $X$；
\item 若 $X,Y$ 独立，则 $H(X\mid Y)=H(X)$，侧信息没有帮助。
\end{itemize}


\subsection{常见错误}


\begin{itemize}
\item 以为“解码器有 $Y$”就可以把 $X$ 压到 $0$ 比特；还需要 $H(X\mid Y)$。
\item 把 Slepian–Wolf 的单源界写成 $L_1\ge H(X)$；有侧信息时是条件熵。
\item 忘记三个下界要同时满足；只满足 $L_1\ge H(X\mid Y)$ 不一定够，还要看总和。
\item 把“随机 binning”当成实际可实现的低复杂度算法；它主要是存在性证明，实际实现需要额外设计。
\end{itemize}


\section*{一页式总览}
\begin{itemize}
\item 无损变长码：$H(X)\le E[L]<H(X)+1$。
\item Kraft：前缀码存在 $\iff\sum_i2^{-l_i}\le1$。
\item Huffman：每次合并两个最小概率，得到最优前缀码。
\item 典型集：$|T_\varepsilon^{(n)}|\le2^{n(H+\varepsilon)}$，概率 $\to1$。
\item 定长几乎无损：$L=n(H+\varepsilon)$，错误概率指数小。
\item 侧信息：$L_1>H(X\mid Y)$。
\item Slepian--Wolf：$L_1\ge H(X\mid Y)$，$L_2\ge H(Y\mid X)$，$L_1+L_2\ge H(X,Y)$。
\end{itemize}

\section*{参考资料}
\begin{itemize}
\item 课堂手写稿（2026-11-19，三页）。
\item 课堂转写（\texttt{11月19日\_转写文本}、\texttt{11月19日 2\_转写文本}）。
\item \textbf{[DS]} 邓小铁、王畅。《离散数学与结构》压缩、熵、Huffman 与信息论相关章节。
\end{itemize}
\end{document}
